\documentclass[final]{beamer}  
  \mode<presentation> {
    \usetheme{Antibes}
    \usecolortheme{orchid}
  }
  \usepackage[utf8]{inputenc}
  \usepackage[spanish]{babel}
  \usepackage{amsmath,amsthm, amssymb, latexsym}
  \usepackage{graphicx}
  \usepackage{lipsum}

  \setbeamertemplate{blocks}[rounded][shadow=true]
  \setbeamerfont{block title}{size=\Large}
  \setbeamerfont{title}{size=\Huge}
  \setbeamerfont{author}{size=\Large}

  \title[\huge Identificación del Sistema Depredador-Presa]{%
    \huge Identificación del Sistema Depredador-Presa\\mediante Redes Neuronales y ANFIS}
  \author{Emmanuel Guerrero Piza \\ Kevin Andrés Forero Guaitero}
  \institute[]{Universidad Distrital Francisco José de Caldas\\Bogotá, Colombia}

\begin{document}

\begin{frame}\

\begin{columns}
    \begin{column}{.47\textwidth}

    \begin{block}{Introducción}
    El modelo de Lotka-Volterra describe la interacción entre presas ($x$) y depredadores ($y$):
    
    $\dot{x} = a x - b x y$ \\[0.3cm]
    $\dot{y} = c x y - d y$
    
    Se discretiza con Euler (T=0.1) y se identifican modelos predictivos usando:
    \begin{itemize}
    \item Red Neuronal (MLP)
    \item ANFIS (RBF + Ridge)
    \end{itemize}
    \end{block}

    \vspace{0.5cm}

    \begin{block}{Metodología}
    \begin{itemize}
    \item 75 trayectorias × 300 pasos = 22,440 datos
    \item 80\% entrenamiento, 20\% prueba
    \item Datos augmentation: 1500 trayectorias diversas
    \item Semilla fija (42) para reproducibilidad
    \end{itemize}
    \end{block}

    \vspace{0.5cm}

    \begin{block}{Arquitectura NN}
    \begin{itemize}
    \item MLP sklearn: (64, 64, 32) neuronas
    \item Activación ReLU
    \item Adam optimizer + early stopping
    \item Data augmentation con ruido N(0,0.3)
    \end{itemize}
    \end{block}

    \vspace{0.5cm}

    \begin{block}{Arquitectura ANFIS}
    \begin{itemize}
    \item 9 funciones de membresía gaussianas por entrada
    \item $\sigma=18$, rejilla uniforme [10,90] → 81 reglas
    \item Regresión Ridge ($\alpha=0.30$)
    \item Entrenamiento instantáneo
    \end{itemize}
    \end{block}

    \end{column}

    \begin{column}{.47\textwidth}

    \begin{block}{Resultados - 1-paso}
    \begin{center}
    \begin{tabular}{lccc}
    \toprule
    Modelo & MAE & MSE & RMSE \\
    \midrule
    NN    & 0.60 & 1.64 & 1.28 \\
    ANFIS & 0.90 & 3.95 & 1.99 \\
    \bottomrule
    \end{tabular}
    \end{center}
    \end{block}

    \vspace{0.5cm}

    \begin{block}{Trayectorias (50 pasos)}
    \begin{center}
    \begin{tabular}{lcc}
    \toprule
    CI & NN & ANFIS \\
    \midrule
    (40,50) & 1.35 & 9.46 \\
    (30,60) & 1.50 & 2.23 \\
    (20,80) & 1.38 & 1.67 \\
    (50,30) & 1.41 & 2.43 \\
    \midrule
    Promedio & 1.41 & 3.95 \\
    \bottomrule
    \end{tabular}
    \end{center}
    \end{block}

    \vspace{0.5cm}

    \begin{alertblock}{Conclusiones}
    \begin{itemize}
    \item La NN logra excelente rendimiento en trayectorias (MAE ~1.4)
    \item ANFIS más rápido pero menor precisión
    \item Los datos augmentados de trayectorias diversas son clave
    \item Semilla fija garantiza reproducibilidad
    \end{itemize}
    \end{alertblock}

    \vspace{0.5cm}

    \begin{block}{Referencias}
    \begin{enumerate}
    \item A.J. Lotka, Elements of physical biology, 1925
    \item V. Volterra, Variazioni e fluttuazioni, 1926
    \item J.-S.R. Jang, ANFIS, IEEE Trans., 1993
    \end{enumerate}
    \end{block}

    \end{column}
\end{columns}

\end{frame}
\end{document}